\chapter{Reading Iceberg Efficiently}
\label{ch:iceberg-read}

\begin{chapterintro}
Before the Iceberg reader opens a single data file, it has already skipped most
of them --- by partition, by manifest, by column statistics, and optionally by
bloom filter. Each layer needs the query's filter in a form it understands, and
the scan report tells you how well each one worked. This chapter covers the
pruning layers, the metrics modes that feed them, and how to read a scan report
as a diagnostic.
\end{chapterintro}

\section{Manifest pruning and metadata filtering}
\tbd

\section{Column statistics and predicate pushdown into Iceberg}
\tbd

\section{Metrics modes: \texttt{full}, \texttt{truncate}, \texttt{counts}, \texttt{none}}
\tbd

\section{Bloom filters}
\tbd

\section{Vectorized reads}
\tbd

\section{Split planning}
\tbd

\section{The metadata tables as a diagnostic}
\tbd

\section{Summary}
\tbd

\exercisehead
\begin{exercises}
\item Run a selective query and, from the scan report, compute the fraction of
      data pruned at each layer (partition, manifest, file).
\end{exercises}
